% FILE: 01_research_question.tex
% PROJECT: crosswalks
% ROLE: type-law-crosswalk-between-standards-and-execution

\section{Research Question}
\label{sec:crosswalks-rq}

\subsection{Problem Statement}
\label{sec:crosswalks-rq-problem}

Process mining tools routinely translate process artifacts between formalisms: from \textsc{BPMN}
to Petri nets for conformance checking, from \textsc{OCEL} to \textsc{XES} for trace-level
analysis, from process trees to WF-nets for alignment computation. Each such translation can
destroy structural information. The problem is not that information is destroyed --- loss is
sometimes mathematically unavoidable --- but that the destruction is \emph{silent}.

\textsc{PM4Py}'s \texttt{convert\_to\_petri\_net(bpmn)} silently approximates
\textsc{OR}-gateways as \textsc{XOR} or \textsc{AND}, drops data objects, annotations, and
lane structure, and returns a \texttt{(PetriNet, Marking, Marking)} triple with no record of
what was lost (source: \texttt{BPMN\_TO\_WFNET.md}). Its \texttt{convert\_to\_event\_log(ocel)}
picks a default case notion without declaration, drops all non-chosen object types, and produces
a \textsc{XES} log with no indication that it was projected from a multi-object \textsc{OCEL}
source (source: \texttt{OCEL\_TO\_XES.md}). The caller has no structural basis for knowing
whether the translated artifact faithfully represents its source.

The research question is: \emph{can every cross-formalism projection in the process intelligence
stack be given a formal type-law contract that names each category of loss, requires an explicit
policy decision before loss occurs, and refuses the projection when loss is not permitted?}

\subsection{Old-Regime Assumption Challenged}
\label{sec:crosswalks-rq-oldregime}

The old-regime assumption is that format conversion is a utility function: it accepts an artifact
in one format and returns an artifact in another format. Loss is considered an implementation
detail, not a first-class structural event. This assumption is instantiated in \textsc{PM4Py}'s
design: conversion functions return converted objects, not conversion verdicts. There is no
type-level distinction between a lossless conversion and an approximating one.

This assumption is challenged on two grounds. First, it is epistemically dishonest: a conformance
score computed on an approximated Petri net is not a conformance score for the original \textsc{BPMN}
process. The score measures conformance to the approximation, not to the model the analyst intended
to check. Second, it is architecturally unsafe for an evidence-grounded system: if the translation
layer can destroy information silently, then the downstream evidence chain --- receipt, checkpoint,
conformance verdict --- is built on an unaudited foundation.

\subsection{Hypothesis}
\label{sec:crosswalks-rq-hypothesis}

\textbf{[AUTHOR THESIS]:} Every projection between the process intelligence formalisms
(\textsc{BPMN}, \textsc{OCEL}, \textsc{XES}, Declare, WF-net, \textsc{POWL}, Process Tree) can
be characterised by a typed \texttt{LossPolicy} contract. This contract requires the caller to
declare --- before the projection executes --- whether loss is refused, permitted with a named
projection, or permitted with a full \texttt{LossReport}. When the policy is \texttt{RefuseLoss}
and loss is structurally inherent (as in \textsc{OCEL} $\to$ \textsc{XES}), the projection
returns a typed \texttt{Refusal} rather than a silently approximated artifact. When the policy
is \texttt{AllowLossWithReport}, the \texttt{LossReport<From, To, DroppedItems>} becomes a
first-class companion to the converted artifact, carrying the audit trail that \textsc{PM4Py}
discards.

The eight crosswalk documents define this contract for each projection direction. The contracts
are research specifications; their enforcement in \texttt{wasm4pm-compat} is the downstream
authorization subject.

\subsection{Success Criteria}
\label{sec:crosswalks-rq-success}

The crosswalks project succeeds when:

\begin{enumerate}
  \item Every major projection direction in the stack has a named \texttt{LossPolicy} contract
        that specifies whether it is universally lossless, conditionally lossless, or inherently
        lossy.
  \item Every category of information destroyed in a lossy projection has a named type
        (e.g.\ \texttt{BpmnToWfNetRefusal::InclusiveGatewayNotRepresentable},
        \texttt{XesExportRefusal::LossRefused}).
  \item The lossless direction (Process Tree $\to$ WF-net) is proven lossless by citation to
        the block-structure theorem (van der Aalst \& Weijters, 2004; Leemans, 2022) and
        enforced at compile time by \texttt{TypedLoopNode<ARITY>} with
        \texttt{Require<\{ ARITY == 2 \}>: IsTrue}.
  \item The \textsc{ALIVE\_001} gate criterion 9 (\texttt{crosswalks/} file count $\geq 3$) is
        met --- verified as MET with actual count 8 (source: \texttt{ALIVE\_GATE\_ASSESSMENT.md}).
\end{enumerate}

The crosswalk status is currently \textbf{PARTIAL}: the structural contracts are defined, but no
explicit \textsc{BLAKE3}-signed receipt chain exists for the directory, and no compilation
evidence confirms that the defined Rust structs are implemented in the actual
\texttt{wasm4pm-compat} crate (source: \texttt{project\_manifest.yaml},
\texttt{receipt\_present: false}).
